\ansclass{12}{Cloud-Native Platforms}
\ans{1} (a) \final{Factor VI, stateless processes}: state in one replica's memory is invisible to the others and dies with the pod. The hold must live in a shared store: \final{PostgreSQL with an expiry column} (or Redis with a TTL). (b) $0.75 \times 1$ GiB $=$ \final{768 MiB} of heap, leaving room for threads, metaspace and buffers. (c) $20 \times 800$ MB $=$ \final{16 GB} against $20 \times 250$ MB $=$ \final{5 GB}: smaller images mean faster scale-out exactly when it matters. (d) A tag can be moved to another image; \final{a digest names exactly the bytes that were tested}, so staging and production run the same artifact.
\careful{``it worked with one replica''. Scaling out is the moment hidden state appears.}
\ans{2} (a) Startup: $5 \times 12 =$ \final{60 s}; out of the Service after $5 \times 3 \approx$ \final{15 s}; restarted after $10 \times 3 \approx$ \final{30 s}. (b) 25\% of 4 $=$ 1 surge pod: \final{at most 5 pods, always 4 ready}. (c) All pods fail readiness at the same time and \final{the whole \texttt{catalog} leaves the Service for 20 s}: a database blip becomes a total outage (in liveness it would even restart the fleet). Readiness checks only the pod's own ability to serve. (d) A CPU limit is enforced per 100 ms period: 20 busy threads use the quota in the first 25 ms and then \final{stall for 75 ms}, so p99 grows while the average CPU looks fine. Set the CPU request, no CPU limit, and always a memory limit.
\careful{putting downstream dependencies in probes. Probes decide about \emph{this} pod; a shared dependency makes every pod fail together.}
\ans{3} (a) $\lceil 4 \times 90/60 \rceil =$ \final{6}; $\lceil 5 \times 96/60 \rceil =$ \final{8}. (b) $4 \to 8 \to 16 \to 32 \to 40$: \final{4 steps, $\approx$ 60 s}, plus 20--40 s of JVM warm-up per pod and minutes for new nodes: \final{2--5 minutes} after the spike started. (c) \final{Pre-scale}: at 09:45 a scheduled job (or a KEDA cron trigger) raises \texttt{minReplicas} to the on-sale size; keep spare nodes (balloon pods); put the waiting room in front of \texttt{inventory}; scale back after the sale. (d) $300 \times 20 = 6\,000$ req/s $/ 150 =$ \final{40 replicas: exactly the maximum}, with no room for losing a node. Raise the maximum (and the node pool), or take load off with the CDN and the cache.
\careful{``the HPA will handle the on-sale''. A reactive autoscaler arrives minutes after a spike that lasts seconds.}
\ans{4} (a) 5\% of 400 $=$ \final{20 req/s}, 25\% $=$ \final{100 req/s}. At 5\%: 835/20 $\approx$ \final{0.7 min}, 2\,312/20 $\approx$ \final{1.9 min}, 18\,527/20 $\approx$ \final{15.4 min}; at 25\%: 0.1, 0.4 and 3.1 min. (b) $30 \times 20 =$ \final{600 requests}: fewer than any of the three needs, so \final{``nothing found'' means nothing}. (c) 09:55 is \final{inside the freeze} (from 09:30): either the expedited canary at 09:42 or wait for the end of the sale. A 5-minute analysis at 5\% sees $5 \times 60 \times 20 = 6\,000$ requests: it detects \final{1\% and 0.5\%, not 0.2\%}. (d) \final{Add a nullable column; write both and read the new one with a fallback; backfill; drop the old column in a later release.} Rolling back the pods does not un-write data: every intermediate state must work with the old and the new version.
\careful{judging a canary by wall-clock time instead of by the number of requests it saw.}
\ans{5} A good drawing: DNS $\to$ CDN $\to$ Ingress (or a load balancer) $\to$ services on Kubernetes ($\to$ PostgreSQL, Redis, Kafka) and the payment provider; thumbnails as functions writing to the object store. Delivery: CI $\to$ registry (signed image) $\to$ commit of the digest in the environment repository $\to$ Argo CD $\to$ cluster $\to$ canary. (a) \final{Services: CaaS (managed Kubernetes); PostgreSQL, Redis, Kafka: DBaaS/managed; thumbnails and webhooks: FaaS; payment and e-mail: SaaS}. (b) \final{In Git}, in the overlay of the environment (the cluster pulls it). Self-heal \final{reverts the manual edit} to what Git says: changes go through a commit (and a revert is the rollback).
\careful{a pipeline that rebuilds the image for each environment: what was tested is not what runs.}
\ans{6} (a) \final{73; 841; 374.50} ($2 \times 0.225 \times 730 + 400 \times 0.115$); \final{188.55} ($3 \times 0.045 \times 730 + 2\,000 \times 0.045$); \final{1\,700; 500} (USD/month). (b) $24 \times 0.192 \times 8$ h $\approx$ \final{US\$ 37 per month}; keeping the 24 extra nodes always on: $24 \times 0.192 \times 730 \approx$ \final{US\$ 3\,364}. You pay for capacity $\times$ time: the peak is cheap, the baseline is not. (c) \final{CDN egress, US\$ 1\,700 ($\approx$ 34\%)}; $\approx 5\,042 / 400\,000 \approx$ \final{US\$ 0.013 per ticket}. Data services ($\approx$ US\$ 1\,470) also cost more than the nodes.
\careful{estimating compute only. Egress, logs, managed data stores and hourly items (NAT, load balancer, control plane) are most of the bill.}
\ans{7} (a) $0.20 \times 10^{-6} + 0.1 \times 0.5 \times 0.0000166667 + 1.00 \times 10^{-6} \approx$ \final{US\$ $2.03 \times 10^{-6}$ per request}; floor: $2 \times 18.02 + 22.27 =$ \final{US\$ 58.31/month}. (b) $58.31 / (2.03 \times 10^{-6} \times 2\,628\,000) \approx$ \final{11 req/s on average} (the class says ``roughly ten''). At 120 req/s the containers need $\lceil 120/50 \rceil = 3$ tasks, \final{US\$ 76}, against \final{$\approx$ US\$ 640} as functions. (c) $2 \times 0.20 + 2 \times 10^6 \times 0.5 \times 1 \times 0.0000166667 \approx$ \final{US\$ 17/month}. (d) \final{Thumbnails and webhooks: functions} (spiky, short, low average, scale from zero). \final{Purchase path: containers}, because cold starts hit p99 exactly at the peak and a burst of functions can exhaust the database connections; with a high steady load they are also cheaper.
\careful{comparing at the peak rate. The bill depends on the \emph{average}; the containers must be sized for the peak (or autoscale).}
